% GENERATED FILE. Edit canonical lessons, then run npm run generate:books.
% canonical-source-hash: fnv1a64:4c1ec8eae820dbf5
% canonical-lessons: ZH-C106-chezhan, ZH-W106-zhan, ZH-C106-ditie, ZH-C106-gongjiaoche, ZH-C106-zixingche, ZH-C106-qi

\chapter{Stops and Stations}
\label{ch:zh-stops-and-stations}
\hlchaptermodality{106}

\begin{chapteropening}
\textbf{By the end of this chapter:} \emph{I can ask for the station, and talk about the metro, buses and bikes.}

Complete the last lesson of chapter 106: I can ask for the station, and talk about the metro, buses and bikes.
\end{chapteropening}

\section[\texorpdfstring{\zh{车站} (chēzhàn)}{车站 (chēzhàn)}]{\zh{车站} (chēzhàn) — station; stop}

\label{lesson:ZH-C106-chezhan}



\begin{quote}
Before the new one: say the Mandarin for soap, then the Mandarin for mirror.
\end{quote}

\subsection*{You'll want to know: \zh{车站}}
\textbf{\zh{车站}} — \emph{chēzhàn} — \textquotedblleft{}station; stop\textquotedblright{}.

\begin{quote}
\textbf{\zh{车站在哪儿}} \emph{chēzhàn zài nǎr} — Where's the station?
\end{quote}

The new character, in pieces: \zh{站} — on the left a five-stroke piece read lì, 'to stand'; on the right a vertical with a short bar to its right, over \zh{口}. The next lesson puts it in your hand.

\begin{cousinweb}
\textbf{\zh{车站}} (\emph{chēzhàn}) is built from \zh{车} 'vehicle' and \zh{站} 'to stand', both inherited from Middle Chinese, \emph{tsyhae} and \emph{treamH} (Guangyun, in Baxter's transcription); the dictionary glosses the second as an everyday word for standing alone. A place where vehicles stand is the usual explanation, but the step is not given there.
\end{cousinweb}

\begin{grammarlens}[title={What you've built}]
One more word for asking where things are.
\end{grammarlens}

\subsection*{Your turn}
\begin{itemize}
\raggedright
  \item \emph{Say it:} \emph{chēzhàn}
  \item \emph{Say it:} \emph{chēzhàn}, once more
  \item \emph{Say it:} \emph{chēzhàn zài nǎr} — Where's the station?
  \item \emph{Recall:} say the Mandarin for soap, then the Mandarin for mirror, then say \emph{chēzhàn} again
\end{itemize}

\begin{tcolorbox}[breakable,colback=teal!4,colframe=teal!35!black,title={Before you move on}]
What does \zh{车站} mean? (\textquotedblleft{}Station; stop\textquotedblright{}.) Say it once more.
\end{tcolorbox}
\section[\texorpdfstring{\zh{站} (zhàn)}{站 (zhàn)}]{\zh{站} — the standing piece, then a mark over \zh{口}}

\label{lesson:ZH-W106-zhan}



\begin{quote}
The word first, then the character inside it.

\begin{itemize}
\raggedright
  \item \emph{Recall:} say \emph{chēzhàn} — \textbf{R1}, one lesson back, before it decays
\end{itemize}

\begin{quote}
\zh{车} \zh{口}
\end{quote}

\emph{Read it:} each one aloud, with nothing said first — \emph{chē}, \emph{kǒu}
\end{quote}

\subsection*{Script — the standing piece, then a mark over \zh{口}}
\hlblockfigure{\detokenize{figures/ZH-W106-zhan-filmstrip.pdf}}{How \zh{站} is written, stroke by stroke}

\begin{quote}
\zh{站}
\end{quote}

\textbf{Two pieces, left before right.} Count the pieces, not the strokes.

\textbf{The standing piece, lì, five strokes}: a dot at the top, a horizontal, a dot on the left and a short left-falling stroke on the right, and a stroke along the bottom that rises to the right.

\textbf{The right piece, five strokes}: a vertical with a short horizontal to its right; then \zh{口} below, as you write it alone: the left side, the top that turns down, the bottom.

5 + 5 = ten.

\subsection*{Your turn}
\begin{itemize}
\raggedright
  \item \emph{Write it:} \zh{站}, all ten, with the model in front of you
  \item \emph{Cover:} write it again from memory, and count the pieces, not the strokes
\end{itemize}

\begin{tcolorbox}[breakable,colback=teal!4,colframe=teal!35!black,title={Before you move on}]
How many pieces in \zh{站}? (\textbf{Two} — the standing piece, then the mark over \zh{口}.) What is the standing piece read? (\emph{lì}.) Which piece you know ends \zh{站}? (\textbf{\zh{口}}.)

Source: \href{https://raw.githubusercontent.com/chanind/hanzi-writer-data/68d10a4b21150cae5e1ebbd223eed289cf32d90c/data/\%E7\%AB\%99.json}{Hanzi Writer Data \zh{站}, pinned stroke paths and medians}.
\end{tcolorbox}
\section[dìtiě]{dìtiě — underground; metro}

\label{lesson:ZH-C106-ditie}



\begin{quote}
Before the new one: say the Mandarin for mirror, then the Mandarin for station.
\end{quote}

\subsection*{You'll want to know: dìtiě}
\emph{dìtiě} — \textquotedblleft{}underground; metro\textquotedblright{}.

\begin{quote}
\emph{wǒmen zuò dìtiě qù} — We'll take the metro.
\end{quote}

Its characters come later; for now it is a word for the ear.

\begin{cousinweb}
\emph{dìtiě} is built from dì 'ground' and tiě 'iron', both inherited from Middle Chinese, \emph{dijH} and \emph{thet} (Guangyun, in Baxter's transcription); the dictionary glosses the second as black metal. The railway under the ground is a modern sense.
\end{cousinweb}

\begin{grammarlens}[title={What you've built}]
One more word for asking where things are.
\end{grammarlens}

\subsection*{Your turn}
\begin{itemize}
\raggedright
  \item \emph{Say it:} \emph{dìtiě}
  \item \emph{Say it:} \emph{dìtiě}, once more
  \item \emph{Say it:} \emph{wǒmen zuò dìtiě qù} — We'll take the metro.
  \item \emph{Recall:} say the Mandarin for mirror, then the Mandarin for station, then say \emph{dìtiě} again
\end{itemize}

\begin{tcolorbox}[breakable,colback=teal!4,colframe=teal!35!black,title={Before you move on}]
What does dìtiě mean? (\textquotedblleft{}Underground; metro\textquotedblright{}.) Say it once more.
\end{tcolorbox}
\section[gōngjiāochē]{gōngjiāochē — city bus}

\label{lesson:ZH-C106-gongjiaoche}



\begin{quote}
Before the new one: say the Mandarin for station, then the Mandarin for metro.
\end{quote}

\subsection*{You'll want to know: gōngjiāochē}
\emph{gōngjiāochē} — \textquotedblleft{}city bus\textquotedblright{}.

\begin{quote}
\emph{gōngjiāochē lái le} — Here comes the bus.
\end{quote}

Its characters come later; for now it is a word for the ear.

\begin{cousinweb}
\emph{gōngjiāochē} is built from gōng 'public', jiāo 'to cross' and \zh{车} 'vehicle', all inherited from Middle Chinese, \emph{kuwng}, \emph{kaew} and \emph{tsyhae} (Guangyun, in Baxter's transcription). The bus is a modern sense of the whole.
\end{cousinweb}

\begin{grammarlens}[title={What you've built}]
One more word for asking where things are.
\end{grammarlens}

\subsection*{Your turn}
\begin{itemize}
\raggedright
  \item \emph{Say it:} \emph{gōngjiāochē}
  \item \emph{Say it:} \emph{gōngjiāochē}, once more
  \item \emph{Say it:} \emph{gōngjiāochē lái le} — Here comes the bus.
  \item \emph{Recall:} say the Mandarin for station, then the Mandarin for metro, then say \emph{gōngjiāochē} again
\end{itemize}

\begin{tcolorbox}[breakable,colback=teal!4,colframe=teal!35!black,title={Before you move on}]
What does gōngjiāochē mean? (\textquotedblleft{}City bus\textquotedblright{}.) Say it once more.
\end{tcolorbox}
\section[zìxíngchē]{zìxíngchē — bicycle}

\label{lesson:ZH-C106-zixingche}



\begin{quote}
Before the new one: say the Mandarin for metro, then the Mandarin for bus.
\end{quote}

\subsection*{You'll want to know: zìxíngchē}
\emph{zìxíngchē} — \textquotedblleft{}bicycle\textquotedblright{}.

\begin{quote}
\emph{zhè shì wǒ de zìxíngchē} — This is my bike.
\end{quote}

Its characters come later; for now it is a word for the ear.

\begin{cousinweb}
\emph{zìxíngchē} is built from zì 'self', xíng 'to go' and \zh{车} 'vehicle', all inherited from Middle Chinese, \emph{dzijH}, \emph{haeng} and \emph{tsyhae} (Guangyun, in Baxter's transcription). The dictionary glosses \emph{haeng} 'to walk' and \emph{dzijH} 'from'; the bicycle is a modern sense of the whole.
\end{cousinweb}

\begin{grammarlens}[title={What you've built}]
One more word for asking where things are.
\end{grammarlens}

\subsection*{Your turn}
\begin{itemize}
\raggedright
  \item \emph{Say it:} \emph{zìxíngchē}
  \item \emph{Say it:} \emph{zìxíngchē}, once more
  \item \emph{Say it:} \emph{zhè shì wǒ de zìxíngchē} — This is my bike.
  \item \emph{Recall:} say the Mandarin for metro, then the Mandarin for bus, then say \emph{zìxíngchē} again
\end{itemize}

\begin{tcolorbox}[breakable,colback=teal!4,colframe=teal!35!black,title={Before you move on}]
What does zìxíngchē mean? (\textquotedblleft{}Bicycle\textquotedblright{}.) Say it once more.
\end{tcolorbox}
\section[qí]{qí — to ride (a bike or a horse)}

\label{lesson:ZH-C106-qi}



\begin{quote}
Before the new one: say the Mandarin for bus, then the Mandarin for bicycle.
\end{quote}

\subsection*{You'll want to know: qí}
\emph{qí} — \textquotedblleft{}to ride (a bike or a horse)\textquotedblright{}.

\begin{quote}
\emph{wǒ qí zìxíngchē qù xuéxiào} — I ride my bike to school.
\end{quote}

Its characters come later; for now it is a word for the ear.

\begin{cousinweb}
\emph{qí} is inherited from Middle Chinese \emph{gje} (Guangyun, in Baxter's transcription), glossed 'to sit astride a horse'.
\end{cousinweb}

\begin{grammarlens}[title={What you've built}]
One more word for asking where things are.
\end{grammarlens}

\subsection*{Your turn}
\begin{itemize}
\raggedright
  \item \emph{Say it:} \emph{qí}
  \item \emph{Say it:} \emph{qí}, once more
  \item \emph{Say it:} \emph{wǒ qí zìxíngchē qù xuéxiào} — I ride my bike to school.
  \item \emph{Recall:} say the Mandarin for bus, then the Mandarin for bicycle, then say \emph{qí} again
\end{itemize}

\begin{tcolorbox}[breakable,colback=teal!4,colframe=teal!35!black,title={Before you move on}]
What does qí mean? (\textquotedblleft{}To ride (a bike or a horse)\textquotedblright{}.) Say it once more.
\end{tcolorbox}
